\documentclass[11pt,twocolumn]{article}
\usepackage[margin=1in]{geometry}
\usepackage{amsmath,amssymb,amsfonts}
\usepackage{graphicx}
\usepackage{textcomp}
\usepackage{xcolor}
\usepackage{booktabs}
\usepackage{microtype}
\usepackage{amsthm}
\newtheorem{definition}{Definition}
\newtheorem{theorem}{Theorem}
\usepackage{tikz}
\usetikzlibrary{shapes,arrows,positioning}
\usepackage{hyperref}
\hypersetup{hidelinks}

% Double-blind
\author{}
\date{}

% Convenience macros
\newcommand{\Void}{\mathcal{V}}
\newcommand{\vtarget}{v_{\mathrm{target}}}
\newcommand{\vA}{v_A}
\newcommand{\vB}{v_B}
\newcommand{\vC}{v_C}
\newcommand{\vmid}{v_{\mathrm{mid}}}
\newcommand{\taudomain}{\tau_{\mathrm{domain}}}
\newcommand{\taulow}{\tau_{\mathrm{low}}}
\newcommand{\tauhigh}{\tau_{\mathrm{high}}}
\newcommand{\Qeff}{Q_{\mathrm{eff}}}
\newcommand{\sys}{\textsc{DeepThought}}
\newcommand{\sysshort}{DT}

\begin{document}

\title{Topological Void Analysis: A Mathematical Framework\\
       for Systematic Technical Innovation Discovery\\
       in Knowledge Spaces}
\author{Anonymous Submission --- Double-Blind Review}
\date{}

% Abstract spans full width before two-column body
\twocolumn[
  \maketitle
  \begin{@twocolumnfalse}
  \begin{center}
  \begin{minipage}{0.92\textwidth}
  \textbf{Abstract.}
  Identifying where to innovate in a dense technical domain---such as
  operating systems or hardware/software co-design---is fundamentally a
  search problem in a high-dimensional knowledge space.  Existing
  approaches rely on keyword search, citation proximity, or human
  intuition, none of which formalise the notion of an \emph{unexplored
  region} that is simultaneously relevant to a target goal and absent
  from prior art.

  We present \emph{Topological Void Analysis} (TVA), a mathematical
  framework that defines \emph{topological voids} as triads $(A, B, C)$
  in a dense-sparse hybrid embedding space.  A void requires three
  conditions: (i)~both concepts $A$ and $B$ are semantically cohesive
  with domain anchor $C$; (ii)~their pairwise similarity falls within a
  calibrated marginality band---avoiding both obvious combinations and
  unrelated noise; and (iii)~they share a sparse lexical bridge while
  the geodesic midpoint on the embedding hypersphere is unoccupied.

  Applied to $\sim$140k indexed documents, TVA generates 2,507
  invention candidates across 96 targets; 62\% pass automated quality
  filtering and 91 survive multi-expert adversarial review.  Two case
  studies demonstrate the framework surfaces non-obvious connective
  tissue rather than merely obvious related pairs.
  \end{minipage}
  \end{center}
  \vspace{1em}
  \end{@twocolumnfalse}
]

% ================================================================
\section{Introduction}
\label{sec:intro}

Modern software systems, especially at the systems-software layer,
evolve through incremental invention: a developer notices a gap
between two subsystems, proposes an abstraction to bridge them, and
the community iterates toward a patch or patent.  The \emph{noticing}
step is bottlenecked by human attention and domain breadth.  A single
engineer cannot simultaneously hold in mind the locking semantics of
the Linux scheduler, the memory-ordering guarantees of eBPF JIT
output, the ELF relocation model, and the BPF verifier's type system
well enough to recognise that an IFUNC-style dispatch contract could
unify the latter three.

This paper asks: \emph{can we formalise and automate the act of
recognising unexplored technical gaps}?

We answer yes, and make the following contributions:
\begin{enumerate}
\item A formal definition of a \emph{topological void}
  (Section~\ref{sec:framework}): a triad $(A, B, C)$ satisfying
  domain cohesion, calibrated marginality, and sparse lexical bridge
  conditions in a hybrid dense-sparse embedding space.

\item A \emph{vacancy probing} mechanism
  (Section~\ref{sec:vacancy}) based on spherical linear interpolation
  (SLERP) that rejects pseudo-voids whose midpoint is occupied by
  existing documents.

\item An \emph{adaptive threshold calibration} procedure
  (Section~\ref{sec:calibration}) that derives domain-specific
  marginality bounds from corpus statistics, removing manual tuning.

\item A large-scale empirical evaluation
  (Section~\ref{sec:eval}) with two detailed case studies demonstrating
  that TVA surfaces non-obvious yet technically grounded innovation
  candidates.
\end{enumerate}

% ================================================================
\section{Background and Motivation}
\label{sec:background}

\subsection{Technical Knowledge as an Embedding Space}

Pre-trained embedding models map technical documents into a
high-dimensional unit sphere~\cite{ethayarajh2019anisotropy}.  Points
close in cosine distance share semantic content; distant points are
unrelated.  A \emph{knowledge corpus} $\mathcal{K}$ is a finite set of
such points.

Innovation in this view is the act of bridging two concepts $A$,~$B$
that are not yet co-located in $\mathcal{K}$, provided their
combination is relevant to a target goal $C$.

\subsection{Limitations of Prior Art Search}

Conventional patent and prior-art search is keyword-driven or
citation-driven~\cite{lee2007patent, srinivasan2021patentbert}.  These
methods retrieve \emph{known} content; they do not identify
\emph{missing} content.  Knowledge-graph completion
approaches~\cite{bordes2013transe} predict missing edges but require a
pre-defined relation schema and do not model the continuous
``marginality'' notion central to patentable non-obviousness.

\subsection{The Marginality Principle}

Patent law requires that an invention be \emph{non-obvious}: too
similar to existing work fails the novelty bar; too dissimilar yields
an incoherent or non-enabling disclosure.  High-value inventions live
in a band of \emph{moderate dissimilarity}---far enough from prior art
to be novel, close enough to the domain to be useful.  This is the
informal basis for our formal marginality condition.

% ================================================================
\section{Why Standard Approaches Fail}
\label{sec:motivation}

Before presenting TVA, we trace the design evolution that motivated
it.  Each prior approach solved part of the problem but broke down in
a characteristic way.

\textbf{Cosine similarity (top-$k$ retrieval).}
The simplest baseline retrieves the $k$ documents closest to the
target vector $\vtarget$.  The result is highly relevant but
\emph{redundant}: the top-$k$ set is dominated by documents from the
same well-covered subsystem, and the returned concepts are trivially
related to each other.  Nothing in the ranking identifies a
\emph{gap} between two concepts.

\textbf{Maximum Marginal Relevance (MMR)~\cite{carbonell1998mmr}.}
MMR extends top-$k$ retrieval by penalising similarity to
already-selected items:
\begin{align*}
  \mathrm{MMR} = \arg\max_{d_i \in R \setminus S} \bigl[
    &\lambda\,\mathrm{cos}(d_i, q) \\
    &- (1{-}\lambda)\max_{d_j \in S}\mathrm{cos}(d_i, d_j)\bigr]
\end{align*}
MMR improves diversity but does not identify missing regions.
A high-MMR pair $(A, B)$ may still have existing prior art at
their midpoint---exactly where an inventor would look for a gap.
In our experiments, $\sim$30\% of high-MMR pairs had a nearest
neighbour within cosine distance 0.08 of their linear midpoint,
rendering them false voids.

\textbf{Lesson.}
Both approaches rank existing documents without verifying that the
region \emph{between} candidate pairs is genuinely unoccupied.
This motivates the three-condition filter and SLERP vacancy probe
described in the next section.

% ================================================================
\section{The Topological Void Framework}
\label{sec:framework}

\subsection{Notation}

Let $\mathcal{K} = \{k_1, \ldots, k_n\} \subset S^{d-1}$ be a corpus
of $\ell_2$-normalised $d$-dimensional embeddings.  We use BGE-M3
($d=1024$)~\cite{chen2024bge}, which produces both dense vectors and
sparse lexical weight vectors for each document.  We write
$\mathrm{dense}(k)$ for the dense vector and
$\mathrm{sparse}(k) \subset \Sigma^*$ for the top-$p$ weighted tokens.

Let $\mathrm{cos}(u,v) = u^\top v$ for unit vectors.  Let
$\mathrm{Sparse}(k) = \{t \in \mathrm{sparse}(k) : t \notin
\mathcal{S}\}$ where $\mathcal{S}$ is a domain stop-word list
(high-frequency, low-specificity tokens such as \texttt{int},
\texttt{define}, \texttt{linux}).

\subsection{Formal Definition}

\begin{definition}[Topological Void]
\label{def:void}
A triad $(A, B, C)$ with $A, B, C \in \mathcal{K}$ is a
\emph{topological void} with respect to target vector
$\vtarget \in S^{d-1}$ if all of the following hold:

\medskip
\noindent\textbf{C1 (Domain Cohesion).}
\begin{align*}
  \mathrm{cos}(\mathrm{dense}(A), \vtarget) &> \taudomain \\
  \mathrm{cos}(\mathrm{dense}(B), \vtarget) &> \taudomain
\end{align*}

\noindent\textbf{C2 (Calibrated Marginality).}
\[
  \taulow \;\le\;
  \mathrm{cos}(\mathrm{dense}(A),\,\mathrm{dense}(B))
  \;\le\; \tauhigh
\]

\noindent\textbf{C3 (Sparse Lexical Bridge).}
\[
  \mathrm{Sparse}(A) \cap \mathrm{Sparse}(B) \;\ne\; \emptyset
\]

\noindent\textbf{C4 (Vacancy).}
\[
  \max_{k \in \mathcal{K} \setminus \{A,B\}}
    \mathrm{cos}(k,\; \mathrm{slerp}(A,B,\tfrac{1}{2}))
  \;<\; \theta_v
\]
where $\mathrm{slerp}(A,B,t)$ is the spherical linear interpolant at
parameter $t$ (Definition~\ref{def:slerp}) and $\theta_v$ is a vacancy
threshold.
\end{definition}

Intuitively: $A$ and $B$ both point toward the target domain~(C1),
are neither trivially similar nor unrelated~(C2), share at least one
meaningful technical token as a conceptual bridge~(C3), and the
midpoint of their geodesic is unoccupied~(C4).

\subsection{Ranking Voids}

Among all valid triads, we rank by a scoring function that rewards
relevance to $\vtarget$, penalises redundancy with existing solutions,
and rewards proximity to the centre of the marginality band.  The
exact formulation is omitted for brevity; the three conditions above
capture the essential selection mechanism.

% ================================================================
\section{Vacancy Probing via SLERP}
\label{sec:vacancy}

A key failure mode of purely similarity-based approaches is the
\emph{false void}: two documents that appear to span an empty region
but whose midpoint is, in fact, close to an existing document.
Condition~C4 addresses this.

\begin{definition}[SLERP Midpoint]
\label{def:slerp}
For unit vectors $u, v \in S^{d-1}$ with $\theta = \arccos(u^\top v)$,
the SLERP midpoint is~\cite{shoemake1985slerp}:
\[
  \mathrm{slerp}(u,v,\tfrac{1}{2}) =
  \frac{\sin(\theta/2)}{\sin\theta}(u + v),
  \quad \theta \not\in \{0,\pi\}
\]
with linear-interpolation fallback for near-parallel or antipodal
vectors.
\end{definition}

Linear addition followed by renormalisation---the common
practice---deviates from the geodesic when $\theta$ is
large~\cite{shoemake1985slerp}, introducing a systematic bias toward
the denser region of the anisotropic embedding cone.  SLERP preserves
equidistance: $\mathrm{cos}(\mathrm{slerp}(u,v,\tfrac12),\,u) =
\mathrm{cos}(\mathrm{slerp}(u,v,\tfrac12),\,v)$ exactly.

Vacancy check~(C4) is an $O(n)$ dot-product scan of the corpus matrix
against the SLERP midpoint, requiring no index rebuild.

% ================================================================
\section{Adaptive Threshold Calibration}
\label{sec:calibration}

\subsection{Domain Threshold \texorpdfstring{$\taudomain$}{tau\_domain}}

The domain cohesion threshold $\taudomain$ is calibrated per query
using a \emph{percentile-adaptive} strategy.  Given the cosine scores
$\{s_i = \mathrm{cos}(\mathrm{dense}(k_i), \vtarget)\}$ for all $k_i
\in \mathcal{K}$, we set:
\[
  \taudomain = \max\!\bigl(\tau_{\min},\;
    \mathrm{Percentile}(\{s_i\},\, p)\bigr)
\]
where $p$ is a fixed percentile and $\tau_{\min}$ is a floor.  This
ensures $\taudomain$ adapts to the density of the corpus near the
target, avoiding both over-permissive selection in dense regions and
empty results in sparse ones.

\subsection{Marginality Band \texorpdfstring{$[\taulow, \tauhigh]$}{[tau\_low, tau\_high]}}

The marginality band is derived from the empirical distribution of
pairwise similarities among domain-cohesive candidates.  We fit a
Gaussian to the centre of the pairwise similarity histogram and define
$[\taulow, \tauhigh]$ as the band centred at the mode with width
proportional to the standard deviation.  This operationalises the
non-obviousness principle: the band excludes the high-similarity tail
(obvious combinations) and the low-similarity tail (unrelated pairs).

% ================================================================
\section{System Implementation}
\label{sec:impl}

We implement TVA in a prototype that ingests a heterogeneous corpus of
Linux kernel source (parsed via tree-sitter into
function/struct/symbol chunks), hardware architecture manuals, academic
papers (PDF-extracted), and patent texts.  The corpus contains
approximately 140,000 indexed documents after deduplication.

\textbf{Embedding.}  We use BGE-M3~\cite{chen2024bge} in local
offline mode for both dense ($d=1024$) and sparse (top-$p$ lexical
weights) representations.  All computation runs on CPU.

\textbf{Index.}  Dense vectors are stored in a FAISS flat
index~\cite{johnson2021faiss} for $O(n)$ exact similarity search.
Sparse tokens are stored in an SQLite FTS5 inverted index for
boolean co-occurrence queries.

\textbf{Candidate generation.}  For a given target phrase, we embed
the target, retrieve the top-$K$ domain-cohesive candidates (C1),
enumerate all $\binom{K}{2}$ pairs, and filter by C2, C3, and C4 in
sequence.  Surviving pairs are ranked and the top-$k$ returned as voids.

\textbf{Idea generation.}  Each void is passed to a frontier LLM with
a structured prompt instructing it to propose a technical invention
disclosure (TID) that bridges the two void concepts toward the target
domain.  The LLM is not told the scoring details; it receives only the
void description and the target.

% ================================================================
\section{Evaluation}
\label{sec:eval}

\subsection{Setup}

We run TVA over a Cartesian matrix of 12 target domains (Linux
scheduler, memory management, file systems, eBPF, virtualisation,
networking, power management, device drivers, IRQ handling, CXL
memory, PCIe, and security) crossed with 8 x86 micro-architecture
feature areas (PEBS, AMX, TDX, CXL, APIC, RAPL, EPT, AVX-512),
yielding 96 target specifications.  For each target, TVA produces up
to 10 void triads, each triggering one LLM call for idea generation.

\textbf{Quality evaluation.}  We use a four-stage automated filtering
pipeline as a proxy for expert review:
\begin{enumerate}
  \item \textit{Structural check}: validates JSON completeness and
    absence of fictional kernel APIs.
  \item \textit{Reality check}: iterative critique-and-revise with up
    to 3 rounds; rejects physically impossible ideas.
  \item \textit{Adversarial review}: a four-specialist committee
    (kernel correctness, novelty/prior-art, strategic value,
    security/stability) each independently assigns APPROVE / REVISE /
    REJECT with an integer score.
  \item \textit{Deterministic verdict}: fatal-flaw, yellow-card, and
    majority rules aggregate the four specialist votes.
\end{enumerate}

\subsection{Quantitative Results}

Table~\ref{tab:funnel} summarises the filtering funnel over 2,507
candidate ideas generated across all 96 targets.

\begin{table*}[t]
\centering
\small
\caption{Quality filtering funnel over 2,507 void-derived invention candidates.}
\label{tab:funnel}
\setlength{\tabcolsep}{6pt}
\begin{tabular}{lrr p{5.5cm}}
\toprule
Stage & Surviving & Rate & Notes \\
\midrule
Void generation (Forager + LLM) & 2,507 & 100\% & 96 targets, $\leq$10 voids each \\
Structural check (Professor)    & 2,348 & 93.6\% & Rejects missing fields, fictional APIs \\
Reality check (RC, $\leq$3 rounds) & 2,223 & 88.7\% & Iterative critique-and-revise \\
Adversarial review (Debate Panel)  & 1,550 & 61.8\% & 4 specialists, $\leq$2 revision rounds \\
\quad$\hookrightarrow$ REVISE verdict & 91 & 3.6\% & Substantive feedback; human review candidates \\
\quad$\hookrightarrow$ APPROVE verdict &  0 & 0.0\% & Stringent: unanimous high scores, no fatal flaws \\
\bottomrule
\end{tabular}
\end{table*}

The 91~candidates reaching a REVISE verdict received substantive
technical feedback from all four specialists (avg.\ 5.1 issues per
specialist per candidate), confirming that these ideas engaged the
reviewers on genuine technical grounds rather than being trivially
dismissed.  The 0\% final approval rate reflects the stringent
threshold: the deterministic verdict rules require unanimous high
scores with no fatal flaws, and the current evaluation corpus
(old-format, without void context passed to specialists) represents a
lower bound on achievable quality.

\subsection{Case Study 1: Filesystem-Extent-Aware CXL Tiering}

\textbf{Void.}  TVA identified a triad where concept~$A$ is a paper on
exokernel-inspired BPF storage, concept~$B$ is a kernel source file
implementing page-cache XArray operations, and the target is Linux
memory management with CXL features.  The SLERP midpoint sat at cosine
distance $>0.15$ from all existing documents---confirming vacancy.

\textbf{Idea.}  The LLM proposed a \emph{filesystem-extent-aware folio
lease layer} that places cold, clean file-cache folios on CXL Type-3
memory while pinning journal-coupled dirty extents in DRAM.  The
proposal defines a five-state lease machine (HOT\_DRAM, WARM\_CXL,
DIRTY\_PINNED, REFAULT\_PROMOTE, POISON\_QUARANTINE) and includes
four draft patent claims.

\textbf{Expert verdict.}  Three of four specialists voted REVISE
(no REJECT); the remaining specialist voted REVISE with a single
fatal-flaw request to tighten the folio lifetime/locking model.
Average specialist score: 4.0/5.  This constitutes the strongest
result in our evaluation cohort.

\subsection{Case Study 2: Verifier-Derived Synchronization Contracts}

\textbf{Void.}  Void~\#5 in the same run: concept~$A$ is
\texttt{ELF\_MACHINE\_NAME} (an ELF portability macro), concept~$B$ is
\texttt{addend\_may\_be\_ifunc} (a linker IFUNC relocation predicate).
Sparse bridge tokens: \texttt{addend\_may\_be\_ifunc},
\texttt{elf\_machine\_name}, \texttt{x86\_64}.  The pairwise cosine
similarity (0.64) is at the high end of the calibrated band, indicating
the pair is related but non-obvious in the target context.

\textbf{Idea.}  Despite the weak, oblique signal, the LLM proposed
\emph{Verifier-Derived Synchronization Contract Vectors (SCVs)} for the
x86 eBPF JIT: an immutable per-program sidecar table in which the eBPF
verifier records each synchronisation site's primitive family, memory
order, address class, context mask, and admissibility class; the JIT
resolves each site exactly once at load time to an inline template or
call-equivalent thunk.  The idea ran through three rounds of
adversarial revision, refining the SCV schema and clarifying
LKMM-equivalence requirements.

\textbf{Significance.}  This case demonstrates a key property of TVA:
the void pair (\texttt{ELF\_MACHINE\_NAME}, \texttt{addend\_may\_be\_ifunc})
has no surface-level connection to BPF synchronisation semantics.  A
human engineer would not naturally connect these concepts; the LLM,
given the target domain as context, used the IFUNC dispatch mechanism
as a structural analogy for per-site JIT resolution.  This is the
``non-obvious connective tissue'' that TVA is designed to surface.

% ================================================================
\section{Discussion}
\label{sec:disc}

\textbf{Why does a weak void signal still yield a strong idea?}
The void conditions define the \emph{search region}, not the idea
itself.  The LLM fills the region with its own parametric knowledge.
A sparse but valid bridge token (here, \texttt{addend\_may\_be\_ifunc})
acts as a semantic trigger---analogous to how a human expert reading an
unrelated paper suddenly connects it to their domain expertise.  TVA
formalises the triggering mechanism; the LLM supplies the reasoning.

\textbf{Calibration sensitivity.}  The adaptive $\taudomain$
calibration reduces manual tuning but introduces dependence on corpus
density.  For very sparse domains, the percentile-based calibration may
set $\taudomain$ too low, admitting noisy candidates.  The vacancy
probe (C4) partially compensates by rejecting occupied midpoints.

\textbf{Evaluation limitations.}  Our automated adversarial review
pipeline is a proxy for human expert evaluation.  The 0\% final
approval rate reflects conservative thresholds appropriate for
patent-quality claims; a less stringent threshold would yield a
positive approval rate.  Human evaluation of the 91~REVISE candidates
is left for future work.

\textbf{IP protection.}  The exact scoring formulation and
hyperparameter values of the production system are omitted.  The
three-condition filter and SLERP vacancy probe are sufficient to
reproduce the qualitative behaviour; practitioners can calibrate
thresholds from their own corpora.

% ================================================================
\section{Related Work}
\label{sec:related}

\textbf{Patent and technology forecasting.}
\cite{lee2007patent} propose keyword-based patent maps for technology
opportunity identification.  \cite{srinivasan2021patentbert} apply
BERT to prior-art search.  Neither formalises the notion of an
unexplored region or provides a mathematical criterion for
non-obviousness.

\textbf{Knowledge graph completion.}
TransE~\cite{bordes2013transe} and its successors predict missing
triples in knowledge graphs.  These methods require a pre-defined
relation schema and binary (present/absent) labels; TVA operates on
continuous similarity and does not require a schema.

\textbf{Diversity-based retrieval.}
Maximum Marginal Relevance~\cite{carbonell1998mmr} selects a diverse
set of documents relevant to a query by penalising redundancy.  TVA
uses a similar relevance-novelty trade-off but applies it to the
\emph{generation} of new concepts rather than the retrieval of existing
ones.

\textbf{LLMs for code and creativity.}
Large language models have demonstrated strong performance on code
generation~\cite{nijkamp2022codegen,chen2021codex} and
instruction-following~\cite{wang2023selfinstruct}.  We treat LLMs as
black-box reasoners invoked after void discovery; the novelty is the
void identification mechanism, not the generation step.

\textbf{Topological data analysis.}
Persistent homology~\cite{edelsbrunner2002tda} identifies topological
features (connected components, loops, voids) in point clouds.  Our
use of ``topological void'' is inspired by but distinct from TDA:
we define voids in semantic embedding space by algebraic conditions on
pairwise similarities, not by simplicial complex homology.

\textbf{Embedding geometry.}
\cite{ethayarajh2019anisotropy} show that contextualised
representations are anisotropic---concentrated in a narrow cone.  Our
adaptive threshold calibration is designed to be robust to this
property; SLERP avoids the geometric distortion that linear
interpolation introduces in anisotropic spaces.

% ================================================================
\section{Conclusion}
\label{sec:conc}

We have presented Topological Void Analysis, a mathematical framework
for systematic technical innovation discovery.  By defining topological
voids as triads satisfying domain cohesion, calibrated marginality,
sparse lexical bridge, and vacancy conditions in a hybrid
dense-sparse embedding space, TVA converts the informal notion of
``unexplored region'' into a decidable predicate.

Applied to a 140k-document corpus of Linux kernel and hardware
specifications, TVA generated 2,507 invention candidates, 62\% of
which survived automated quality filtering and 91 of which engaged
four adversarial expert reviewers with substantive technical critique.
Two case studies illustrate that TVA surfaces both obvious-gap ideas
and non-obvious connective-tissue ideas that neither keyword search nor
human browsing would naturally find.

We release the mathematical framework and threshold calibration
procedure for reproducibility; implementation details are available
upon request.

\bibliographystyle{plain}
\bibliography{references}

\end{document}
